\begin{textsample}{POS Dim 4 – vem\_ed\_35 – Score 13.00 – vem\_ed\_35/t188.txt}  \label{ex:f4_pos_003}
BOAS PRÁTICAS Com sabor de café e chocolate Alessandro José Padin Ferreira Fatec Praia Grande Nos dois primeiros semestres de 2025, vivi uma das experiências mais significativas da minha \textbf{trajetória} docente.
Por meio dos Projetos Colaborativos Internacionais (PCIs), desenvolvidos na abordagem COIL (Collaborative Online International Learning), tive a oportunidade de orientar dois PCIs entre a Fatec Praia Grande e o Departamento Universitario Obrero y Campesino (DUOC), do Chile.
Foram encontros virtuais que se transformaram em \textbf{espaços} de troca humana, cultural e afetiva.
No primeiro PCI, trabalhei com estudantes do curso de Comércio Exterior.
A proposta \textbf{parecia} objetiva: discutir possibilidades de inserção e fortalecimento do café brasileiro especial no mercado chileno, mas logo percebi que não seria possível falar apenas da mercadoria.
No Brasil, café também carrega memória, identidade e emoção.
Ao longo das atividades, incentivei os alunos a compreenderem que internacionalizar um produto também \textbf{significa} internacionalizar narrativas.
Não bastava apresentar números de exportação ou \textbf{tendências} de mercado.
Era necessário mostrar aos colegas chilenos como esse fruto faz parte da nossa história cotidiana.
Falamos sobre Santos, tão próxima da Fatec Praia Grande, cidade cuja \textbf{trajetória} está profundamente ligada ao porto por onde o produto brasileiro alcançou o mundo.
Conversamos sobre o cheiro do café passado nas manhãs brasileiras, sobre as \textbf{reuniões} em família e o hábito quase ritualístico de oferecer uma xícara a quem chega.
Em muitos momentos, percebi que os estudantes começaram a enxergar o comércio exterior de outra maneira.
Não apenas como negociações internacionais, mas como encontros culturais.
As apresentações produzidas por eles mostravam o crescimento do café brasileiro no Chile e a valorização crescente das alternativas especiais.
Mas o que mais me marcou foi observar como os alunos passaram a construir conexões emocionais em torno do tema.
Uma das equipes trouxe poemas e suas memórias afetivas.
Naquele instante, compreendi que o projeto havia alcançado algo maior do que eu imaginava inicialmente, pois \textbf{deixava} de ser apenas um produto de exportação e se tornava linguagem cultural.
Os estudantes chilenos demonstravam enorme \textbf{curiosidade} ao perceber essa relação afetiva do brasileiro.
Muitos comentavam como o Chile, historicamente ligado ao consumo de chá, \textbf{vinha} desenvolvendo novos hábitos ligados aos cafés especiais e às cafeterias contemporâneas.
As conversas eram espontâneas, e pouco a pouco as barreiras linguísticas iam \textbf{desaparecendo}.
Chocolate orgânico No segundo semestre de 2025, desenvolvi um novo PCI, desta vez com estudantes de Gestão Empresarial.
O tema escolhido foi o chocolate orgânico brasileiro.
Mais uma vez, procurei estimular os alunos a compreenderem que por trás de um produto existem territórios, pessoas, tradições e modos de vida.
As discussões nos levaram ao sul da Bahia, às cooperativas agrícolas, à agricultura familiar e ao sistema cabruca, técnica tradicional de cultivo do cacau sob a sombra da Mata Atlântica.
Falamos sobre sustentabilidade, preservação ambiental e sobre como pequenas comunidades brasileiras transformaram o chocolate orgânico em símbolo de identidade regional.
Foi especialmente interessante observar o envolvimento dos estudantes ao pesquisarem o universo dos chocolates bean-to-bar (\textbf{significa} do “grão à barra” e define a produção de forma artesanal e integral), dos produtos de origem única e das marcas brasileiras que começam a conquistar mercados internacionais a partir de narrativas ligadas à sustentabilidade e à rastreabilidade.
Ao longo dos dois projetos, percebi algo que muitas vezes a rotina acadêmica nos faz esquecer: a educação possui um enorme poder de aproximação humana.
Mesmo separados por milhares de quilômetros, estudantes brasileiros e chilenos criaram vínculos genuínos.
No início, havia insegurança com o idioma, receio de falar espanhol, medo de errar.
Depois, \textbf{surgiram} risadas, espontaneidade e confiança.
Como professor, foi extremamente gratificante acompanhar esse processo.
Vi alunos descobrirem que a internacionalização não pertence apenas às grandes universidades ou aos programas \textbf{presenciais} no exterior.
Pode acontecer, também, dentro da sala de aula virtual, por meio do diálogo, da escuta e da colaboração.
Competências interculturais O reconhecimento institucional do projeto destacou competências interculturais, comunicação, habilidades digitais e trabalho em equipe global.
Mas acredito que houve um aprendizado ainda mais profundo: a compreensão de que produtos carregam histórias e que nenhuma estratégia internacional é realmente eficiente quando ignora a cultura das pessoas envolvidas.
Ao final daqueles encontros, tive a sensação de que café e chocolate eram apenas o ponto de partida.
O que realmente atravessou as fronteiras entre Brasil e Chile foram memórias, identidades e experiências compartilhadas.
E talvez seja essa a principal lembrança que guardarei dos PCIs de 2025: a certeza de que ensinar também é criar pontes humanas, algumas delas feitas de tecnologia, outras de \textbf{palavras}, e muitas delas do aroma inesquecível de um café recém passado.

% matched lemmas: curiosidade, deixar, desaparecer, espaço, palavra, parecer, presencial, reunião, significar, surgir, tendência, trajetória, vir
\end{textsample}
